\subsection{Bounded Diameter Trees}
\label{sec:bounded_diameter}

Results for trees with diameter constraint $D$.

\subsubsection{Diameter 5 - Uniform Distribution}

\begin{table}[H]
\centering
\caption{Trees with $D = 5$, uniform weights}
\label{tab:diam5_uniform}
\begin{tabular}{|c|c|c|c|c|}
\hline
$n$ & Avg Depth & Max Depth & Avg Time (ms) & Std Dev \\
\hline
% Data from: bounded_diameter_5_uniform.csv
\csvreader[head to column names]{data/bounded_diameter_5_uniform.csv}{}{%
\n & \avgdepth & \maxdepth & \avgtime & \stddev \\
}
\hline
\end{tabular}
\end{table}
\begin{figure}[H]
\centering
\begin{tikzpicture}
\begin{axis}[
    xlabel={Tree size ($n$)},
    ylabel={Average depth},
    width=0.48\textwidth,

    grid=major,
    legend pos=north west,
    mark size=2pt
]
\addplot[blue, mark=*, thick] table[x=n, y=avgdepth, col sep=comma] {data/bounded_diameter_5_uniform.csv};
\legend{Avg Depth}
\end{axis}
\end{tikzpicture}
\hfill
\begin{tikzpicture}
\begin{axis}[
    xlabel={Tree size ($n$)},
    ylabel={Runtime (ms)},
    width=0.48\textwidth,

    grid=major,
    legend pos=north west,
    mark size=2pt
]
\addplot[red, mark=square*, thick] table[x=n, y=avgtime, col sep=comma] {data/bounded_diameter_5_uniform.csv};
\legend{Avg Time}
\end{axis}
\end{tikzpicture}
\caption{Performance metrics: Trees with $D = 5$, uniform weights}
\label{fig:diam5_uniform}
\end{figure}


\subsubsection{Diameter 10 - Uniform Distribution}

\begin{table}[H]
\centering
\caption{Trees with $D = 10$, uniform weights}
\label{tab:diam10_uniform}
\begin{tabular}{|c|c|c|c|c|}
\hline
$n$ & Avg Depth & Max Depth & Avg Time (ms) & Std Dev \\
\hline
% Data from: bounded_diameter_10_uniform.csv
\csvreader[head to column names]{data/bounded_diameter_10_uniform.csv}{}{%
\n & \avgdepth & \maxdepth & \avgtime & \stddev \\
}
\hline
\end{tabular}
\end{table}
\begin{figure}[H]
\centering
\begin{tikzpicture}
\begin{axis}[
    xlabel={Tree size ($n$)},
    ylabel={Average depth},
    width=0.48\textwidth,

    grid=major,
    legend pos=north west,
    mark size=2pt
]
\addplot[blue, mark=*, thick] table[x=n, y=avgdepth, col sep=comma] {data/bounded_diameter_10_uniform.csv};
\legend{Avg Depth}
\end{axis}
\end{tikzpicture}
\hfill
\begin{tikzpicture}
\begin{axis}[
    xlabel={Tree size ($n$)},
    ylabel={Runtime (ms)},
    width=0.48\textwidth,

    grid=major,
    legend pos=north west,
    mark size=2pt
]
\addplot[red, mark=square*, thick] table[x=n, y=avgtime, col sep=comma] {data/bounded_diameter_10_uniform.csv};
\legend{Avg Time}
\end{axis}
\end{tikzpicture}
\caption{Performance metrics: Trees with $D = 10$, uniform weights}
\label{fig:diam10_uniform}
\end{figure}


\subsubsection{Diameter 20 - Uniform Distribution}

\begin{table}[H]
\centering
\caption{Trees with $D = 20$, uniform weights}
\label{tab:diam20_uniform}
\begin{tabular}{|c|c|c|c|c|}
\hline
$n$ & Avg Depth & Max Depth & Avg Time (ms) & Std Dev \\
\hline
% Data from: bounded_diameter_20_uniform.csv
\csvreader[head to column names]{data/bounded_diameter_20_uniform.csv}{}{%
\n & \avgdepth & \maxdepth & \avgtime & \stddev \\
}
\hline
\end{tabular}
\end{table}
\begin{figure}[H]
\centering
\begin{tikzpicture}
\begin{axis}[
    xlabel={Tree size ($n$)},
    ylabel={Average depth},
    width=0.48\textwidth,

    grid=major,
    legend pos=north west,
    mark size=2pt
]
\addplot[blue, mark=*, thick] table[x=n, y=avgdepth, col sep=comma] {data/bounded_diameter_20_uniform.csv};
\legend{Avg Depth}
\end{axis}
\end{tikzpicture}
\hfill
\begin{tikzpicture}
\begin{axis}[
    xlabel={Tree size ($n$)},
    ylabel={Runtime (ms)},
    width=0.48\textwidth,

    grid=major,
    legend pos=north west,
    mark size=2pt
]
\addplot[red, mark=square*, thick] table[x=n, y=avgtime, col sep=comma] {data/bounded_diameter_20_uniform.csv};
\legend{Avg Time}
\end{axis}
\end{tikzpicture}
\caption{Performance metrics: Trees with $D = 20$, uniform weights}
\label{fig:diam20_uniform}
\end{figure}


\subsubsection{Impact of Weight Distributions ($D = 10$)}

\textbf{Random Distribution:}

\begin{table}[H]
\centering
\caption{Trees with $D = 10$, random weight distribution}
\label{tab:diam10_random}
\begin{tabular}{|c|c|c|c|c|}
\hline
$n$ & Avg Depth & Max Depth & Avg Time (ms) & Std Dev \\
\hline
% Data from: bounded_diameter_10_random.csv
\csvreader[head to column names]{data/bounded_diameter_10_random.csv}{}{%
\n & \avgdepth & \maxdepth & \avgtime & \stddev \\
}
\hline
\end{tabular}
\end{table}
\begin{figure}[H]
\centering
\begin{tikzpicture}
\begin{axis}[
    xlabel={Tree size ($n$)},
    ylabel={Average depth},
    width=0.48\textwidth,

    grid=major,
    legend pos=north west,
    mark size=2pt
]
\addplot[blue, mark=*, thick] table[x=n, y=avgdepth, col sep=comma] {data/bounded_diameter_10_random.csv};
\legend{Avg Depth}
\end{axis}
\end{tikzpicture}
\hfill
\begin{tikzpicture}
\begin{axis}[
    xlabel={Tree size ($n$)},
    ylabel={Runtime (ms)},
    width=0.48\textwidth,

    grid=major,
    legend pos=north west,
    mark size=2pt
]
\addplot[red, mark=square*, thick] table[x=n, y=avgtime, col sep=comma] {data/bounded_diameter_10_random.csv};
\legend{Avg Time}
\end{axis}
\end{tikzpicture}
\caption{Performance metrics: Trees with $D = 10$, random weight distribution}
\label{fig:diam10_random}
\end{figure}


\textbf{Exponential Distribution:}

\begin{table}[H]
\centering
\caption{Trees with $D = 10$, exponential weight distribution}
\label{tab:diam10_exponential}
\begin{tabular}{|c|c|c|c|c|}
\hline
$n$ & Avg Depth & Max Depth & Avg Time (ms) & Std Dev \\
\hline
% Data from: bounded_diameter_10_exponential.csv
\csvreader[head to column names]{data/bounded_diameter_10_exponential.csv}{}{%
\n & \avgdepth & \maxdepth & \avgtime & \stddev \\
}
\hline
\end{tabular}
\end{table}
\begin{figure}[H]
\centering
\begin{tikzpicture}
\begin{axis}[
    xlabel={Tree size ($n$)},
    ylabel={Average depth},
    width=0.48\textwidth,

    grid=major,
    legend pos=north west,
    mark size=2pt
]
\addplot[blue, mark=*, thick] table[x=n, y=avgdepth, col sep=comma] {data/bounded_diameter_10_exponential.csv};
\legend{Avg Depth}
\end{axis}
\end{tikzpicture}
\hfill
\begin{tikzpicture}
\begin{axis}[
    xlabel={Tree size ($n$)},
    ylabel={Runtime (ms)},
    width=0.48\textwidth,

    grid=major,
    legend pos=north west,
    mark size=2pt
]
\addplot[red, mark=square*, thick] table[x=n, y=avgtime, col sep=comma] {data/bounded_diameter_10_exponential.csv};
\legend{Avg Time}
\end{axis}
\end{tikzpicture}
\caption{Performance metrics: Trees with $D = 10$, exponential weight distribution}
\label{fig:diam10_exponential}
\end{figure}


\subsubsection{Comparison Across Diameter Bounds}

\begin{table}[H]
\centering
\caption{Performance comparison for different diameter bounds ($n = 100$, uniform weights)}
\label{tab:diameter_comparison}
\begin{tabular}{|c|c|c|c|c|}
\hline
$D$ & Avg Depth & Max Depth & Avg Time (ms) & Std Dev \\
\hline
% Data from: bounded_diameter_comparison.csv
\csvreader[head to column names]{data/bounded_diameter_comparison.csv}{}{%
\diameter & \avgdepth & \maxdepth & \avgtime & \stddev \\
}
\hline
\end{tabular}
\end{table}
\begin{figure}[H]
\centering
\begin{tikzpicture}
\begin{axis}[
    xlabel={Tree size ($n$)},
    ylabel={Average depth},
    width=0.48\textwidth,

    grid=major,
    legend pos=north west,
    mark size=2pt
]
\addplot[blue, mark=*, thick] table[x=n, y=avgdepth, col sep=comma] {data/bounded_diameter_comparison.csv};
\legend{Avg Depth}
\end{axis}
\end{tikzpicture}
\hfill
\begin{tikzpicture}
\begin{axis}[
    xlabel={Tree size ($n$)},
    ylabel={Runtime (ms)},
    width=0.48\textwidth,

    grid=major,
    legend pos=north west,
    mark size=2pt
]
\addplot[red, mark=square*, thick] table[x=n, y=avgtime, col sep=comma] {data/bounded_diameter_comparison.csv};
\legend{Avg Time}
\end{axis}
\end{tikzpicture}
\caption{Performance metrics: Performance comparison for different diameter bounds ($n = 100$, uniform weights)}
\label{fig:diameter_comparison}
\end{figure}


\textbf{Observations:}
\begin{itemize}
    \item Diameter constraints naturally limit decision tree depth
    \item Smaller diameter bounds result in more balanced tree structures
    \item Diameter constraint forces star-like structures for small $D$
    \item Weight distribution effects diminish as diameter decreases
    \item Performance remains consistent across diameter variations
    \item Trade-off between tree diameter and search efficiency
\end{itemize}
